\chapter{Stato dell'arte}
\label{cap:stato-arte}

\section{Sicurezza della supply chain del software}
\todo{Pratiche di sviluppo sicuro e riferimenti normativi, ad esempio \cite{nist-ssdf}.}

\section{Software Bill of Materials}
Una \sbom{} è un inventario strutturato dei componenti che compongono un software, con relazioni di
dipendenza, versioni e licenze. I formati più diffusi sono \emph{SPDX}~\cite{spdx} e \emph{CycloneDX}~\cite{cyclonedx}.

\begin{table}[htbp]
  \centering
  \caption{Confronto tra formati \sbom{} (da completare).}
  \label{tab:formati-sbom}
  \begin{tabularx}{\linewidth}{@{}lXX@{}}
    \toprule
    \textbf{Aspetto} & \textbf{SPDX} & \textbf{CycloneDX}\\
    \midrule
    Ambito principale & \todo{} & \todo{}\\
    Formati di serializzazione & \todo{} & \todo{}\\
    Supporto a vulnerabilità & \todo{} & \todo{}\\
    \bottomrule
  \end{tabularx}
\end{table}

\section{Certificati di sicurezza e firma del software}
\todo{Certificati X.509, firma del codice, attestazioni di provenienza e infrastrutture a chiave pubblica.}

\section{Blockchain per l'integrità e la tracciabilità}
\todo{Registro distribuito, smart contract, ancoraggio di hash~\cite{nakamoto}.}
